% !TEX root = main.tex
% This file contains the "LLM-as-a-Judge Assessment" subsection,
% converted from Experiments.md.
% Required packages (add to your preamble if not already loaded):
%   \usepackage{booktabs}
%   \usepackage{graphicx}
%   \usepackage{amsmath}
%   \usepackage{amssymb}
%   \usepackage{xcolor}
%   \usepackage{listings}  % only if you want fancy code formatting

\section{Experiments}
\label{sec:experiments}

We evaluate the proposed privacy-preserving query rewriter on two
complementary axes. \emph{Privacy} measures how much
task-non-essential sensitive content survives in the rewritten query
that is exposed to the remote LLM, while \emph{utility} measures
whether the rewrite is still informative enough for a frozen local
aggregator (or a remote model) to produce a high-quality response.
Section~\ref{sec:exp-setup} describes the models, metrics, and
baselines; Section~\ref{sec:exp-main} reports the privacy--utility
trade-off across three local aggregators of different scales; and
Section~\ref{sec:llm-judge} presents a fine-grained LLM-as-a-Judge
analysis of failure modes.

\subsection{Experimental Setup}
\label{sec:exp-setup}

\paragraph{Models.}
We instantiate our rewriter $\pi_\theta$ from
\textsc{Qwen-2.5-3B-Instruct}, and train it with GRPO under the reward
defined in \S\ref{sec:reward}. A \emph{single} rewriter checkpoint is
shared across all evaluation settings: we deliberately do \emph{not}
re-train one rewriter per downstream aggregator, so as to test whether
the learned rewriting behaviour generalises beyond a particular
backbone. The remote answerer used by the \textit{Direct LLM} ceiling
is \textsc{GPT-4o} accessed through an internal proxy, and is also
used as the LLM-as-a-Judge in
Section~\ref{sec:llm-judge}.\footnote{For the fine-grained
flag-and-score judge in \S\ref{sec:llm-judge} we use the more
cost-efficient \textsc{GPT-4o-mini}; for the pairwise utility judge we
use full \textsc{GPT-4o} with temperature $0$.} On the trusted side,
we evaluate three frozen local aggregators of increasing scale:
\textsc{Qwen-2.5-1.5B-Instruct}, \textsc{Qwen-2.5-7B-Instruct}, and
\textsc{Llama-3.1-8B-Instruct}. These aggregators are kept frozen and
are \emph{not} part of the optimisation loop: at inference time, the
rewritten query $y$ is sent to the remote LLM, and the retrieved
answer is recombined with the original query $x$ on the local
aggregator to produce the final response. Training and decoding
hyper-parameters are reported in Appendix~\ref{app:training}.

\paragraph{Evaluation data.}
We evaluate on the benchmark introduced in \S\ref{sec:data}, which
provides every test instance with an explicit
$(\mathcal{E},\mathcal{N})$ partition of essential and non-essential
sensitive spans. The benchmark consists of two complementary subsets.
The \textbf{General} split aggregates the \emph{Common} and
\emph{Wild} portions of our reverse-synthesised general-domain set and
covers 11 privacy-relevant domains and 20 task types, capturing the
typical case in which non-essential disclosures are incidental to the
user's request. The \textbf{Medical} split is the board-exam-style
challenge set designed to stress the over-removal failure mode, in
which clinical signals are simultaneously sensitive and
answer-decisive. Reporting the two splits separately allows us to
distinguish a method's behaviour on \emph{redundant} versus
\emph{answer-bearing} sensitive content. Per-split statistics are
summarised in Table~\ref{tab:data-stats} (Appendix~\ref{app:data}).

\paragraph{Evaluation metrics.}
We adopt complementary metrics for the two axes.

\textit{Privacy.} For each rewrite $y$, an item $c\in\mathcal{N}$ is
flagged as leaked if either its token-level Jaccard coverage in $y$
exceeds $\tau_{\text{cov}}$ \emph{or} its rescaled BERTScore
precision~\citep{zhang2020bertscore} against $y$ exceeds
$\tau_{\text{bert}}$, i.e., we capture both verbatim and paraphrastic
disclosure. Items already mentioned in $\mathcal{E}$ are excluded
from leakage counting. Based on this primitive we report three
sample-level rates:
\textbf{SLR} (\emph{sensitive-leak rate}), the fraction of samples on
which any non-essential sensitive span is leaked, our headline
privacy score;
\textbf{NE} (\emph{non-essential leak rate}), the per-item leakage
rate over all spans in $\mathcal{N}$, which controls for samples
carrying many sensitive items; and
\textbf{EA} (\emph{essential-attribute retention}), the per-item
retention rate over $\mathcal{E}$, which complements SLR by measuring
how much answer-bearing context survives the rewrite. Lower SLR/NE
and higher EA indicate better privacy-aware behaviour. To facilitate
comparison with prior PII-redaction work, we also track a \textbf{PII}
leak rate restricted to the canonical PII categories produced by the
detector used in~\citet{prediax} on the same rewrites; this metric is
reported in Appendix~\ref{app:pii}.

\textit{Utility.} Because the General and Medical splits have
qualitatively different answer formats, we evaluate task utility with
two judging protocols, both implemented with \textsc{GPT-4o}. On
\emph{General}, we follow the pairwise-with-swap protocol of
\citet{papillon}: given the original question, the local pipeline's
answer $C_L$ (driven by the rewrite $y$) is compared against the
remote ceiling answer $C_T$ produced by \textsc{GPT-4o} on the
unprotected query. We query the judge in both orderings and take the
swap-consistent decision; the reported \textit{Util.} score is the
fraction of samples on which $C_L$ is judged at least as good as
$C_T$. On \emph{Medical}, where each item has a verified
multiple-choice answer label, we instead use a strict single-choice
consistency judge that returns $1$ iff the option selected in $C_L$
exactly matches the gold label. The full prompts and parsing rules
are given in Appendix~\ref{app:judge-utility}.

We additionally report two summary numbers per aggregator block:
the \textbf{average utility} across General and Medical, and the
gap \textbf{$\Delta$ vs.\ Local}, defined as the difference between
this average and the corresponding \textit{Local Only} utility under
the same aggregator. $\Delta$ vs.\ Local isolates the \emph{net}
contribution of off-device assistance after the privacy filter has
been applied, which is the quantity an end-user actually cares about.

\paragraph{Baselines.}
We compare with two representative families of query-side privacy
defences, plus two reference systems that bound the
privacy--utility trade-off from above and below.
\textbf{Direct LLM} forwards the unprotected query to the remote
\textsc{GPT-4o} and serves as the \emph{utility ceiling}; it is
unprotected and is not a comparable system on the privacy axis.
\textbf{Local Only} runs the original query entirely on the frozen
local aggregator and exposes nothing off-device, serving as the
\emph{privacy reference}: privacy metrics are not applicable, and its
utility is the floor that any off-device-assisted method must
strictly beat to be useful. \textbf{Prediax}~\citep{prediax} is a
rule-based PII redactor that masks canonical PII categories
(e.g.\ \texttt{<PERSON>}, \texttt{<DATE\_TIME>}) with placeholders;
its rewrite is independent of the local aggregator, so we report it
once. \textbf{Papillon}~\citep{papillon} is the strongest prior
LLM-based rewriter we are aware of; it ties its rewriter to the local
aggregator backbone, so we instantiate one Papillon per aggregator
(\textsc{Qwen-2.5-1.5B/7B} and \textsc{Llama-3.1-8B}) using the
official prompts. For all systems we use greedy decoding for the
rewriter and identical decoding settings for the downstream
answerer. Implementation details and prompt templates of the

\subsection{Main Results on Privacy-Preserving Query Rewriting}
\label{sec:exp-main}

\paragraph{Overall comparison.}\textit{Local Only} under the same aggregator.

\paragraph{Overall comparison.} Table~\ref{tab:main} evaluates the privacy--utility trade-off of different query rewriting strategies. The local aggregator is used only as the downstream executor for measuring task utility, while the privacy metrics are computed on the rewritten queries. Since the Qwen-2.5-1.5B results are not yet complete, we focus on the completed Qwen-2.5-7B and Llama-3.1-8B settings.

Overall, our method achieves the best average utility among protected methods in both completed settings. With Qwen-2.5-7B as the downstream aggregator, our method obtains an average utility of 78.7, improving over Local Only by 10.1 percentage points and outperforming Papillon by 2.4 points. With Llama-3.1-8B, our method achieves an average utility of 80.5, exceeding Local Only by 6.6 points and Papillon by 9.8 points. These results indicate that our rewriting strategy better preserves task-relevant information while providing privacy protection.

The improvement is particularly clear on the Medical split. For Qwen-2.5-7B, our method improves Medical utility from 47.3 to 62.0 over Local Only, while reducing leakage compared with Papillon. A similar trend is observed with Llama-3.1-8B, where our method reaches 66.7 Medical utility, compared with 60.7 for Local Only and 49.2 for Papillon. This suggests that our method is especially beneficial for domain-specific queries, where preserving useful semantic information under privacy constraints is more challenging.

On the General split, Papillon achieves slightly lower leakage on some metrics, such as SLR under the Llama-3.1-8B setting. However, our method consistently provides higher utility and maintains low leakage levels, leading to a stronger overall privacy--utility balance. Importantly, the same Qwen-2.5-3B rewriter is used across all evaluation settings, suggesting that the rewriting behavior is not tied to a specific downstream aggregator.

\subsection{LLM-as-a-Judge Assessment}
\label{sec:llm-judge}

\paragraph{Protocol.}
To complement the automatic privacy--utility metrics, we conduct a
fine-grained LLM-as-a-Judge evaluation (GPT-4o-mini) on 500 rewrites per
system (1{,}500 in total) for our method and two strong baselines,
\emph{Papillon} and \emph{Prediax}. Given the original question and the
rewritten query, the judge produces (i) flags for four privacy-leak
categories and five utility-issue categories, (ii) any newly
hallucinated sensitive attributes, and (iii) a re-identification risk
score $r\in[0,1]$, a utility score $u\in[0,1]$. The four leak
and five utility categories are defined in
Appendix~\ref{app:judge-rubric}.

\paragraph{Aggregate results.}
We observe a clear negative correlation between $r$ and $u$
(Pearson $-0.43$ overall), confirming the expected privacy--utility
tension. Within this tension, our rewrites occupy the most balanced
region of the score space: $11.2\%$ of our samples fall into the
\emph{ideal} high-utility / low-risk quadrant ($u>0.7$ and $r\le 0.5$)
--- comparable to Prediax ($14.8\%$) --- while our \emph{worst-case}
quadrant (low utility and high risk) contains only $3.6\%$ of samples,
less than half of Prediax's $9.0\%$. This suggests that our method
produces fewer catastrophic rewrites than Prediax while remaining more
informative than Papillon.

\begin{figure}[t]
  \centering
  \includegraphics[width=1\linewidth]{figures/08_quadrant_distribution.png}
  \caption{Joint distribution of utility score $u$ and re-identification
  risk score $r$ across the three systems. Our method occupies the
  most balanced region, with relatively few samples in the
  low-utility / high-risk quadrant.}
  \label{fig:quadrant-distribution}
\end{figure}

\paragraph{Privacy-failure mode analysis.}
Breaking the privacy-issue rate down by leak type
(Table~\ref{tab:leaks}) shows that the three systems fail in
qualitatively different ways. Prediax relies heavily on placeholder
substitution (e.g., \texttt{<PERSON>}, \texttt{<DATE\_TIME>}) and
consequently exhibits the highest \texttt{direct\_leak} rate
($51.8\%$ of samples), because non-name tokens such as identifiers,
diagnoses, or device names are left untouched. Papillon performs the
most aggressive paraphrasing and therefore reduces direct leakage to
$15.8\%$, but its rewrites retain enough correlated context to trigger
\texttt{combination\_leak} in $50.8\%$ of samples. Our method sits
between these two regimes, with the lowest \emph{severe} leakage rate
(direct $+$ semantic combined: $25.2\%$, vs.\ $20.0\%$ / $54.4\%$ for
Papillon / Prediax) and the lowest fraction of multi-issue privacy
failures ($4.2\%$ vs.\ $7.2\%$ / $12.2\%$).

\begin{table*}[t]
  \centering
  \small
  \caption{Privacy-leak rates (\%) by category. Bold marks the best
  (lowest, except for ``\,$\ge 1$ issue'' / ``\,$\ge 2$ issues'' columns
  where lower is also better) value per column among comparable systems.}
  \label{tab:leaks}
  \begin{tabular}{lcccccc}
    \toprule
    System   & direct        & semantic     & inferential  & combination   & $\ge 1$ issue   & $\ge 2$ issues \\
    \midrule
    Papillon & 15.8          & 4.2          & 12.8         & \textbf{50.8} & 76.4            & 7.2            \\
    Prediax  & 51.8          & 2.6          & 10.6         & 34.6          & 86.2            & 12.2           \\
    \textbf{Ours} & 23.0     & \textbf{2.2} & \textbf{8.0} & 39.6          & \textbf{68.4}   & \textbf{4.2}   \\
    \bottomrule
  \end{tabular}
\end{table*}

\paragraph{Utility-failure mode analysis.}
On the utility side, the dominant failure across all systems is
\texttt{lost\_constraint}, which is consistent with the observation
that anonymization tends to drop execution-relevant specifics. Severe
utility failures (\texttt{task\_drift}, \texttt{over\_generalization})
remain rare ($<0.3\%$) but, when they occur, drive the utility score
down to ${\sim}0.20$. Our method exhibits slightly fewer multi-issue
utility failures ($2.2\%$) than the two baselines ($6.8\%$ / $5.4\%$),
indicating that the constraint loss it does incur is typically
confined to a single mild weakening rather than compounding into task
drift.

\paragraph{Hallucinated sensitive content.}
A non-trivial $7.2\%$ of all rewrites introduce sensitive attributes
that are not entailed by the original question
($108 / 1{,}500$ samples, $130$ spans), spanning categories such as
contact identifiers ($11.5\%$ of spans), locations, demographics, and
family or medical context. The rate is comparable across systems
(Papillon $9.2\%$, Ours $6.6\%$, Prediax $5.8\%$), suggesting this is
a general hazard of LLM-driven rewriting and motivates explicit
constraints against attribute hallucination in future work.


% ... existing code ...
\newcommand{\best}[1]{\textbf{#1}}                  % best in column
\newcommand{\sbest}[1]{\underline{#1}}              % second best
\newcommand{\aggrhead}[2]{%
  \rowcolor{#1}\multicolumn{12}{c}{\textbf{#2}}\\}
\newcommand{\todofill}{{\color{gray!70}?}}          % placeholder for 1.5B
\definecolor{ourrow}{HTML}{FFF6D6}                  % light yellow for "Ours"
\definecolor{qwenblock}{HTML}{FFF1E6}               % light orange block header
\definecolor{llamablock}{HTML}{EAF7EA}              % light green block header
\definecolor{metricgray}{HTML}{F5F5F5}              % light gray header
\definecolor{refblock}{HTML}{EEF2FB}                % light blue for reference

% =============================================================================
%  TABLE 1 (MAIN). Privacy-utility comparison across aggregator scales.
% =============================================================================
\begin{table*}[t]
\centering
\setlength{\tabcolsep}{6pt}
\renewcommand{\arraystretch}{1.22}
\caption{%
  \textbf{Privacy--utility comparison of query rewriting strategies under three
  frozen local aggregators.} \best{Bold} and \sbest{underline} mark the best
  and second-best \emph{protected} results within each aggregator block.
} 
\label{tab:main}
\small
\resizebox{\textwidth}{!}{%
\begin{tabular}{l l c c c c c c c c c c}
\toprule
\textbf{Methods} & \textbf{Rewriter}
& \multicolumn{4}{c}{\textbf{General}}
& \multicolumn{4}{c}{\textbf{Medical}}
& \textbf{Avg.\,Util.}\,$\uparrow$
& \textbf{$\Delta$ vs.\,Local}\,$\uparrow$ \\
\cmidrule(lr){3-6}\cmidrule(lr){7-10}
& &
SLR\,$\downarrow$ &
NE\,$\downarrow$ &
EA\,$\uparrow$ &
Util.\,$\uparrow$ &
SLR\,$\downarrow$ &
NE\,$\downarrow$ &
EA\,$\uparrow$ &
Util.\,$\uparrow$ & & \\
\midrule

% --- Reference / aggregator-agnostic baselines -------------------------------
Direct LLM & raw query
& \multicolumn{3}{c}{\textit{unprotected}} & 100.0$^{\star}$
& \multicolumn{3}{c}{\textit{unprotected}} & 79.5
& 89.8 & --- \\
Prediax & rule-based PII
& 89.6 & 53.7 & 96.2 & 89.2
& 1.8 & 0.6 & 95.3 & 64.2
& 76.7 & --- \\

% --- Qwen-1.5B aggregator (TBD) ----------------------------------------------
\midrule
\aggrhead{qwenblock}{Local aggregator: \textit{Qwen-2.5-1.5B-Instruct}}
Local Only & original query
& \multicolumn{3}{c}{\textit{on-device}} & \todofill
& \multicolumn{3}{c}{\textit{on-device}} & \todofill
& \todofill & 0.0 \\
Papillon & Qwen-2.5-1.5B
& \todofill & \todofill & \todofill & \todofill
& \todofill & \todofill & \todofill & \todofill
& \todofill & \todofill \\
\rowcolor{ourrow}
\textbf{Ours} & Qwen-2.5-3B (\emph{shared})
& 8.0 & 2.2 & 59.8 & \todofill
& 4.2 & 1.5 & 71.0 & \todofill
& \todofill & \todofill \\

% --- Qwen-7B aggregator ------------------------------------------------------
\midrule
\aggrhead{qwenblock}{Local aggregator: \textit{Qwen-2.5-7B-Instruct}}
Local Only & original query
& \multicolumn{3}{c}{\textit{on-device}} & 89.9
& \multicolumn{3}{c}{\textit{on-device}} & 47.3
& 68.6 & 0.0 \\
Papillon & Qwen-2.5-7B
& \best{6.9} & \best{2.0} & \sbest{63.4} & \sbest{93.5}
& 55.0 & 23.9 & \best{80.6} & 59.0
& 76.3 & \sbest{$+$7.7} \\
\rowcolor{ourrow}
\textbf{Ours} & Qwen-2.5-3B (\emph{shared})
& \sbest{8.0} & \sbest{2.2} & \best{59.8} & \best{95.3}
& \best{4.2} & \best{1.5} & \sbest{71.0} & \best{62.0}
& \best{78.7} & \best{$+$10.1} \\

% --- Llama-3.1-8B aggregator -------------------------------------------------
\midrule
\aggrhead{llamablock}{Local aggregator: \textit{Llama-3.1-8B-Instruct}}
Local Only & original query
& \multicolumn{3}{c}{\textit{on-device}} & 87.2
& \multicolumn{3}{c}{\textit{on-device}} & 60.7
& 74.0 & 0.0 \\
Papillon & Llama-3.1-8B
& \best{3.0} & \sbest{4.6} & \sbest{52.4} & \sbest{92.2}
& 46.7 & 16.7 & 61.7 & 49.2
& 70.7 & \sbest{$-$3.3} \\
\rowcolor{ourrow}
\textbf{Ours} & Qwen-2.5-3B (\emph{shared})
& \sbest{8.0} & \best{2.2} & \best{59.8} & \best{94.3}
& \best{4.2} & \best{1.5} & \best{71.0} & \best{66.7}
& \best{80.5} & \best{$+$6.6} \\
\bottomrule
\end{tabular}%
}

\vspace{2pt}
{\footnotesize
$^{\star}$\,Utility ceiling, not a protected method.\quad
``\textit{on-device}'' indicates that no information is exposed off-device,
so privacy metrics are not applicable.\quad
``?'' marks results pending due to ongoing experiments.
}
\end{table*}